\documentclass[conference]{IEEEtran}
\IEEEoverridecommandlockouts

\usepackage{cite}
\usepackage{amsmath,amssymb,amsfonts}
\usepackage{algorithmic}
\usepackage{graphicx}
\usepackage{textcomp}
\usepackage{xcolor}
\usepackage{booktabs}
\usepackage{array}
\usepackage{multirow}
\usepackage{url}
\usepackage{hyperref}
\usepackage{tikz}
\usetikzlibrary{shapes.geometric, arrows.meta, positioning, shadow, calc}

\def\BibTeX{{\rm B\kern-.05em{\sc i\kern-.025em b}\kern-.08em
    T\kern-.1667em\lower.7ex\hbox{E}\kern-.125emX}}

\begin{document}

\title{Predictive Semantic Cache Invalidation for Vector-Based Code Search and RAG Systems\\
\large Review-1 Academic Project Proposal and Assessment Report
}

\author{
    \IEEEauthorblockN{\textbf{Aditya Bhandari}, \textbf{Chebolu Tarun}, \textbf{K. Mohnish}}
    \IEEEauthorblockA{\textit{Department of Computer Science and Engineering} \\
    \textit{School of Computing and Information Technology}\\
    Vellore Institute of Technology, Chennai, India \\
    Email: mohnish.maitreya2022@vitstudent.ac.in}
    \and
    \IEEEauthorblockN{\textbf{Dr. Sivagami M}}
    \IEEEauthorblockA{\textit{Project Guide, Department of Computer Science and Engineering}\\
    \textit{School of Computing and Information Technology}\\
    Vellore Institute of Technology, Chennai, India}
}

\maketitle

\begin{abstract}
Codebase vector search engines and Retrieval-Augmented Generation (RAG) pipelines convert source code functions into dense vector embeddings for semantic navigation, automated code generation, and AI pair-programming. However, software repositories evolve continuously across Git commits. Re-indexing every function upon each commit incurs prohibitive $O(N)$ computational and API token costs. Conversely, naively re-embedding only directly edited files leaves dependent caller functions with stale vector representations, degrading LLM context quality. This paper proposes a native, graph-aware predictive cache invalidation framework. By extracting Abstract Syntax Tree (AST) structures, Git diff statistics, and NetworkX topological features (including PageRank, shortest-path distance, and graph transition descriptors), our pipeline trains a Random Forest regressor to predict embedding drift ($1 - \text{cosine\_similarity}$). A Weighted BFS Decay algorithm propagates invalidation signals along reverse call-graph edges to selectively refresh stale caches. The proposed framework aims to significantly reduce re-embedding overhead while preserving high retrieval recall without requiring heavy external server daemons.
\end{abstract}

\begin{IEEEkeywords}
Vector Databases, Semantic Cache Invalidation, Retrieval-Augmented Generation (RAG), AST Normalization, Call Graph Topology, Machine Learning.
\end{IEEEkeywords}

% ===============================================================
\section{Introduction and Problem Formulation}
% ===============================================================

\subsection{Context and Technical Relevance}
The modern software engineering ecosystem has been fundamentally transformed by AI-assisted code intelligence tools, intelligent IDE assistants (e.g., Cursor, GitHub Copilot, Sourcegraph Cody), and codebase-level Retrieval-Augmented Generation (RAG) platforms. These applications operate by partitioning a software repository into fine-grained code snippets (typically individual functions or class methods) and projecting them into a high-dimensional vector space using neural code embedding models such as UniXcoder, CodeBERT, or sentence-transformers. When a developer submits a natural language query or prompts an AI pair-programmer, the system performs a vector similarity search against a vector database (e.g., ChromaDB, Qdrant, or Pinecone) to retrieve the most semantically relevant code snippets as context for a Large Language Model (LLM).

However, software repositories are dynamic, constantly changing environments. As software development teams continuously merge Git commits, push refactorings, and modify function signatures, the underlying code entities evolve. This evolution introduces a major operational challenge: **vector cache staleness**. If a code entity is edited or its surrounding context changes, its existing cached embedding in the vector database becomes obsolete. Existing industry practices suffer from a strict trade-off between full repository re-indexing ($O(N)$ compute cost) and naive direct diff tracking (\texttt{changed\_only}), which ignores dependent caller functions.

\subsection{Problem Definition and Semantic Drift Taxonomy}
To quantitatively analyze how code updates impact vector indices, we establish a formal taxonomy of code evolution across Git commits ($C_i \to C_{i+1}$):

\subsubsection{Direct Semantic Drift}
Direct Semantic Drift occurs when a function $f$ undergoes internal logic modifications, parameter signature changes, or statement updates directly within the Git diff. In this scenario, the source code text of $f$ itself is modified, causing its vector embedding $v_f$ to drift directly from $v_f^{(i)}$ to $v_f^{(i+1)}$.

\subsubsection{Indirect Semantic Drift}
Indirect Semantic Drift occurs when function $f$ remains unedited in Git, but one of its downstream callees $g$ (where $f$ calls $g$) undergoes logic modifications. Under context-aware chunking regimes (where prompt text incorporates callee interface stubs, docstrings, or MD5 hashes), modifying callee $g$ mutates the contextual prompt text of caller $f$. As a result, caller $f$ undergoes indirect vector drift despite being unedited in Git diffs.

\subsection{Project Objectives}
To address vector cache staleness effectively, the proposed project is defined around seven concrete technical objectives:
\begin{enumerate}
    \item \textbf{Construct a Native AST Code Parsing Engine:} Build a lightweight Python AST parser to extract function and method entities along with symbol definitions.
    \item \textbf{Develop a 25-Signal Topological Feature Engine:} Compute structural, edit volatility, centrality (PageRank, betweenness), shortest-path distance, and McCabe cyclomatic complexity features.
    \item \textbf{Train a Predictive ML Drift Model:} Train a Random Forest regressor on historical Git commit pairs to accurately forecast continuous vector distance drift ($1 - \text{cosine\_similarity}$).
    \item \textbf{Design a Weighted BFS Decay Propagation Algorithm:} Develop a graph cascade algorithm to propagate invalidation signals along reverse call-graph edges to capture indirect caller drift.
    \item \textbf{Implement an AST Cosmetic Normalization Filter:} Perform round-trip AST compilation to strip whitespace, formatting, and docstrings, ignoring non-functional code edits.
    \item \textbf{Deliver a Production Vector Store Integration Plugin:} Package the system into a production-ready extension (\texttt{llama\_index\_integration.py}) for LlamaIndex vector store ingestion pipelines (ChromaDB/Qdrant).
    \item \textbf{Benchmark Against Baseline Strategies:} Systematically evaluate performance against full re-indexing, direct diff tracking (\texttt{changed\_only}), and fixed-hop heuristics across Cost Ratio \%, Recall@K, MAE, and decision latency.
\end{enumerate}

% ===============================================================
\section{Scope, Expected Outcomes, and Feasibility}
% ===============================================================

\subsection{Scope and Project Boundaries}
The scope of this project is strictly defined around function-level static analysis and vector store invalidation.

\subsubsection{In-Scope System Capabilities}
\begin{itemize}
    \item Analysis restricted to Python software repositories.
    \item Fine-grained analysis at the individual function and class method level.
    \item Static call-graph construction and topological feature extraction.
    \item Selective cache invalidation targeting vector databases for code RAG applications.
\end{itemize}

\subsubsection{Explicit System Boundaries and Non-Coverage}
To maintain clear engineering boundaries, the following aspects are explicitly outside the system scope:
\begin{itemize}
    \item \textbf{Dynamic Method Dispatch:} Dynamic runtime polymorphic overrides are resolved statically using module import symbol tables and suffix-matching heuristics rather than dynamic runtime tracing.
    \item \textbf{Non-Python Programming Languages:} The primary implementation targets Python source code; multi-language support (e.g., C++, Java, Rust) is reserved for future Tree-sitter extension.
    \item \textbf{External Network APIs and Runtime Dependencies:} External cloud service calls or non-static data streams are not tracked in static call graphs.
\end{itemize}

\subsection{Expected System Outcomes}
Upon completion, the project will deliver the following key technical outcomes:
\begin{itemize}
    \item \textbf{Predictive Vector Drift Model:} A trained ML regressor predicting continuous cosine embedding distance drift ($1 - \text{cosine\_similarity}$).
    \item \textbf{Call-Graph Decay Propagation:} A functional Weighted BFS Decay algorithm that captures indirect caller drift along reverse call edges.
    \item \textbf{Selective Invalidation Engine:} A high-precision filter that isolates stale code entities while leaving fresh cached vectors intact.
    \item \textbf{Production Integration Plugin:} A production-ready extension (\texttt{llama\_index\_integration.py}) implementing \texttt{CodeGraphNodeParser} and \texttt{PredictiveCacheFilter} for LlamaIndex.
    \item \textbf{Substantial Overhead Reduction:} Target reduction in re-embedding compute and API token costs by up to 90\% compared to full re-indexing.
    \item \textbf{Rigorous Benchmark Evaluation:} Comprehensive empirical evaluation against standard baselines across Recall@K, Cost Ratio \%, Fidelity MAE, and sub-0.15s decision latency.
\end{itemize}

\subsection{Technical and Resource Feasibility Analysis}
The project feasibility is thoroughly evaluated across five primary dimensions:

\subsubsection{Hardware and Software Requirements}
The framework is designed to run on standard commodity hardware (minimum 8 GB RAM, standard CPU/GPU). Software dependencies are lightweight, requiring Python 3.9+, NetworkX, Scikit-Learn, and LlamaIndex.

\subsubsection{Dataset Availability}
Extensive open-source Git repository histories (such as \texttt{psf/black}) provide rich, real-world commit diffs and co-evolution data for model training and evaluation.

\subsubsection{Model Availability}
Pre-trained, state-of-the-art neural code embedding models (\texttt{microsoft/unixcoder-base}, \texttt{sentence-transformers/all-MiniLM-L6-v2}) are publicly accessible for ground-truth vector generation.

\subsubsection{Implementation Feasibility}
By building the core parsing and topological engine using native Python \texttt{ast} and NetworkX graph algorithms, the system operates at a target algorithmic complexity of $O(V + E \log V)$ with a target local memory footprint under $200\text{MB}$ RAM, avoiding heavy external server daemons.

\subsubsection{Deployment Feasibility}
The system packages its output as a LlamaIndex \texttt{TransformComponent}, allowing seamless drop-in integration into existing production vector store pipelines (ChromaDB, Qdrant).

% ===============================================================
\section{Literature Review and Research Gap Synthesis}
% ===============================================================
The challenge of maintaining semantic consistency in vector-based retrieval indices as underlying codebases evolve sits at the intersection of repository mining, change analysis, code embeddings, RAG architectures, and cache invalidation. Below, we review seminal papers categorized into key technical domains.

\subsection{Repository Mining and Structural Change Analysis}

\subsubsection{ChangeDistiller: Fine-Grained AST Tree Differencing}
\emph{Fluri et al. \cite{Fluri2025} (IEEE TSE 2025):} Pioneered the use of Abstract Syntax Tree (AST) differencing to move beyond line-level textual diffs. By extracting tree edit operations (insertions, deletions, moves, and updates), ChangeDistiller enables precise identification of structural code modifications such as method renames and parameter list changes. In our framework, we adopt this tree-edit philosophy to filter out cosmetic edits prior to feature generation.

\subsubsection{Node Feature Enhancement for Code Difference Extraction}
\emph{Zhang et al. \cite{Zhang2024} (IEEE AUTEEE 2024):} Extended AST differencing by introducing a node feature enhancement module that incorporates semantic node attributes. This allows graph-matching algorithms to distinguish genuine node movements from spurious textual edits. We incorporate this concept into our Stage 1 AST normalization engine.

\subsubsection{Logical Coupling Mining from Software Changes}
\emph{Wetzlmaier et al. \cite{Wetzlmaier2014} (IEEE IWSM 2014):} Demonstrated that logical dependencies between heterogeneous artifacts can be mined from co-evolution histories across multi-year development cycles. Their findings show that change coupling provides strong implicit dependency signals, motivating our historical volatility feature set.

\subsection{Dependency Graphs and Commit Semantics}

\subsubsection{SDG Commit Classification via Graph Convolutional Networks}
\emph{Zhang et al. \cite{Zaixing2023} (IEEE QRS 2023):} Constructed System Dependency Graphs (SDGs) from commit diffs and applied program slicing to isolate change impact graphs. A Graph Convolutional Network (GCN) then classifies commits into corrective, perfective, and adaptive categories. This validates that graph-structured dependency representations outperform flat text diffs when modeling commit impact.

\subsubsection{UntCC: Untangling Composite Commits via Structural-Semantic Embeddings}
\emph{Jin et al. \cite{Jin2026} (IEEE TSE 2026):} Addressed composite commit untangling by combining LLaMA-3.2-3B text embeddings with graph autoencoders over code-change graphs. Their work proves that joint semantic-structural representations successfully separate interleaved code modifications across complex Git commits.

\subsubsection{Keyword-Connected PDGs for Fine-Grained Clone Detection}
\emph{Wu et al. \cite{Wu2025} (IEEE TR 2025):} Connected Program Dependency Graph (PDG) nodes through shared keywords to improve fine-grained semantic clone detection beyond what syntactic comparison alone achieves. This demonstrates that linking structural graphs with keyword attributes enhances semantic sensitivity.

\subsubsection{Maintenance Consistency Prediction in Code Clones}
\emph{Zhang et al. \cite{Zhang2020} (IEEE Access 2020):} Demonstrated that machine-learned predictive models operating on structural clone attributes can accurately forecast whether newly introduced code clones require consistent maintenance propagation. This confirms that ML-based propagation prediction is practically valuable.

\subsection{Code Embedding Models and Semantic Drift}

\subsubsection{UniXcoder, CodeBERT, and GraphCodeBERT}
\emph{Guo et al. \cite{guo2022unixcoder}, Feng et al. \cite{feng2020codebert}, Guo et al. \cite{guo2020graphcodebert}:} Established pre-trained transformer architectures for code representations. UniXcoder uses multi-modal mask attention to project code snippets into high-dimensional vector spaces, serving as our primary baseline embedding backbone.

\subsubsection{code2vec and ASTminer for Python Code Embeddings}
\emph{Ngo et al. \cite{Ngo2023} (IEEE SEAI 2023):} Explored code2vec with AST path sampling for generating distributed vector representations of Python code. They demonstrated that AST-path-based representations capture structural code properties effectively across programming languages.

\subsubsection{Longitudinal Code Embedding Evolution and Social Dynamics}
\emph{He et al. \cite{Yulong2026} (Journal of Comput. Sci. 2026):} Modeled codebase evolution by tracking code embedding trajectories longitudinally over open-source repository histories, proving that vector embeddings capture structural evolution and semantic drift over time.

\subsubsection{Memory-Augmented Feature Drift Management}
\emph{Wu et al. \cite{Qilong2026} (Knowledge-Based Systems 2026):} Addressed feature drift in visual place recognition embeddings using an age-weighted memory queue with exponential decay to model descriptor staleness. This concept of controlled freshness management directly mirrors our commit-driven vector drift prediction.

\subsection{Retrieval-Augmented Generation and Cache Invalidation}

\subsubsection{Enterprise RAG Architectures and Index Staleness}
\emph{Lewis et al. \cite{lewis2020retrieval}, Gao et al. \cite{gao2023retrieval}, Ersoy \& Er{\c{s}}ah{\i}n \cite{Ersoy2025}:} Formalized Retrieval-Augmented Generation (RAG) and surveyed enterprise implementations. Ersoy et al. systematically evaluated sparse, dense, and hybrid RAG architectures, confirming that static vector indices suffer measurable staleness costs as source documents evolve.

\subsubsection{Real-Time Event-Driven RAG Index Updates}
\emph{Miji{\'c} \& Isakovi{\'c} \cite{Mijic2026} (IEEE IT 2026):} Designed a production-grade streaming RAG pipeline for CRM applications that reduces document-to-query propagation latency to under 3.1 seconds using event-driven stream processing, underscoring the operational necessity of index freshness.

\subsubsection{Trust-Aware Multi-Agent Knowledge Graphs}
\emph{Essam et al. \cite{Essam2026} (IEEE MECO 2026):} Used confidence-calibrated knowledge graphs as shared semantic memory surfaces for multi-agent software engineering pipelines, showing that threshold gating prevents low-quality upstream decisions from propagating errors downstream.

\subsubsection{CacheSense: Selective Document Cache Invalidation}
\emph{Dang et al. \cite{Dang2026} (IEEE CAIBDA 2026):} Introduced a freshness-aware semantic caching framework for LLM serving backends that selectively invalidates document-level cache entries based on cosine similarity thresholds following knowledge base updates.

\subsubsection{Virtual Data Center Re-Embedding Optimization}
\emph{Satpathy et al. \cite{Satpathy2024} (IEEE TGCN 2024):} Formulated selective re-embedding as a resource optimization problem using genetic algorithms to plan cost-minimizing virtual data center migration under quality constraints.

\subsubsection{Athena: Transformer + PDG Change Impact Analysis}
\emph{Yan et al. \cite{yan2024athena} (ACM FSE 2024):} Fused transformer embeddings with Program Dependence Graphs (PDGs) for method-level change impact set prediction. While Athena targets developer change impact, our system is a **commit-driven proactive vector store cache invalidation engine** that predicts continuous embedding drift ($1 - \text{cosine\_similarity}$) for AST entities and selectively refreshes vector stores prior to retrieval queries.

\subsection{Literature Comparison and Research Gap Synthesis}
Table \ref{tab:lit_review} provides a comparative summary of existing research baselines alongside our proposed solution.

\begin{table*}[htbp]
\caption{Comparative Literature Analysis: Existing Research Baselines vs. Our Proposed System}
\label{tab:lit_review}
\centering
\begin{tabular}{p{3.2cm}p{3.5cm}p{4.2cm}p{5.2cm}}
\toprule
\textbf{Approach / Reference} & \textbf{Core Mechanism} & \textbf{Identified Limitations} & \textbf{Our Proposed Solution} \\
\midrule
Full Re-Indexing \cite{gao2023retrieval} & Re-embeds 100\% of files on every commit. & $O(N)$ token and compute costs; fails to scale to large repos. & **Selective Re-Embedding:** Targets re-embedding to ~1.4\%--5\% of nodes, aiming to save >90\% cost. \\
LlamaIndex \texttt{refresh\_ref\_docs} & File-hash matching at the document level. & Re-embeds entire files if 1 line changes; misses caller drift across files. & **AST Entity Resolution:** Function-level invalidation with cosmetic comment filtering. \\
ChangeDistiller \cite{Fluri2025} & Tree-differencing for AST edit extraction. & Focuses on software history mining, not vector embedding maintenance. & **AST + Drift Predictor:** Uses AST diffs as features for ML vector drift prediction. \\
CacheSense \cite{Dang2026} & Document-level selective cache invalidation. & Operates on generic text documents, lacking call-graph topological signals. & **Graph-Aware Invalidation:** Uses NetworkX PageRank \& BFS Decay over code call graphs. \\
Athena \cite{yan2024athena} & Transformer embeddings + PDG for impact sets. & Targets developer impact set analysis rather than vector store cache freshness. & **Predictive ML Invalidation:** Predicts continuous cosine drift ($1 - \text{sim}$) for RAG indices. \\
\bottomrule
\end{tabular}
\end{table*}

\subsubsection{Limitations of Existing Approaches}
Prior work in vector store maintenance relies on brute-force full re-indexing or flat document-hash matching. Brute-force re-indexing incurs $O(N)$ token cost, while document-hash matching ignores dependency structures, leaving dependent caller functions with stale cached embeddings. Conversely, software change impact tools (e.g., Athena) analyze developer change sets but do not predict vector space drift or interface with production RAG vector stores.

\subsubsection{Synthesized Research Gap}
Existing research lacks a **lightweight, graph-aware predictive invalidation framework** that connects static AST call-graph topology with machine-learned vector drift prediction ($1 - \text{cosine\_similarity}$) specifically designed for production vector database ingestion pipelines.

\subsubsection{How Our Proposed Project Addresses the Gap}
Our project addresses this gap by combining: (1) round-trip AST cosmetic comment filtering, (2) 25 topological NetworkX features, (3) Random Forest ML drift regression, (4) Weighted BFS Decay reverse propagation, and (5) a native LlamaIndex integration plugin (\texttt{PredictiveCacheFilter}).

% ===============================================================
\section{Overview of Feature Taxonomy}
% ===============================================================

Vector search models project source code snippets into a dense vector space $\mathbb{R}^d$. Semantic distance between the original vector $v_f^{(i)}$ and the post-commit vector $v_f^{(i+1)}$ is measured via cosine distance:
\begin{equation}
D_{\text{cosine}}(v_A, v_B) = 1 - \frac{v_A \cdot v_B}{\|v_A\|_2 \|v_B\|_2}
\end{equation}

When code evolves across commits ($C_i \to C_{i+1}$), we formulate cache invalidation as predicting whether the ground-truth embedding distance $D_{\text{cosine}}(v_f^{(i)}, v_f^{(i+1)})$ exceeds a predetermined drift threshold $\tau$.

Our native feature engine extracts 25 structural, historical, and code metrics grouped into five mathematical categories:

\subsection{Git Edit Distance and Volatility (6 Features)}
\begin{itemize}
    \item \texttt{file\_lines\_added}, \texttt{file\_lines\_deleted}: Quantifies the total line-level edit volume within the modified source file.
    \item \texttt{entity\_size}, \texttt{entity\_modification\_size}: Total line count of the target function and the exact number of modified lines within its body.
    \item \texttt{modification\_frequency}, \texttt{previous\_drift}: Historical change frequency of the entity across past commit windows and its prior observed drift.
\end{itemize}

\subsection{Graph Topology and Centrality (5 Features)}
\begin{itemize}
    \item \texttt{in\_degree}, \texttt{out\_degree}: Direct count of caller functions depending on this entity and callee functions invoked by this entity.
    \item \texttt{pagerank}, \texttt{closeness}, \texttt{betweenness}: NetworkX topological centrality metrics measuring structural importance within the repository call graph.
\end{itemize}

\subsection{Change Propagation and Shortest Path (4 Features)}
\begin{itemize}
    \item \texttt{distance\_to\_modified\_directed}, \texttt{distance\_to\_modified\_undirected}: Shortest directed and undirected path distance from the target entity to the nearest edited node in Git.
    \item \texttt{modified\_dependents\_count}, \texttt{modified\_dependencies\_count}: Immediate count of modified caller and callee neighbors.
\end{itemize}

\subsection{Personalized PageRank and Graph Transition Descriptors (7 Features)}
\begin{itemize}
    \item \texttt{pagerank\_impact}: NetworkX Personalized PageRank (PPR) score computed by seeding jump probabilities specifically on Git-modified nodes.
    \item \texttt{gtd\_change\_class}, \texttt{gtd\_local\_edges\_added}, \texttt{gtd\_local\_edges\_removed}, \texttt{gtd\_local\_edge\_churn}, \texttt{gtd\_mean\_drift}, \texttt{gtd\_drift\_variance}: Structural call-graph transition descriptors tracking local edge perturbations.
\end{itemize}

\subsection{Native AST Code Complexity Metrics (3 Features)}
\begin{itemize}
    \item \texttt{cyclomatic\_complexity}: Native McCabe Cyclomatic Complexity computed by counting decision branch AST nodes ($1 + \sum (\text{\texttt{If}}, \text{\texttt{For}}, \text{\texttt{While}}, \text{\texttt{Except}}, \text{\texttt{BoolOp}} )$).
    \item \texttt{ast\_node\_count}: Total count of statement nodes present in the parsed AST.
    \item \texttt{max\_nesting\_depth}: Maximum recursive control structure block depth (\texttt{If}, \texttt{For}, \texttt{While}, \texttt{Try}, \texttt{With}).
\end{itemize}

% ===============================================================
\section{System Methodology and Component Design}
% ===============================================================

The proposed predictive semantic cache invalidation framework operates as a highly optimized, end-to-end processing pipeline divided into three distinct operational phases: (1) Offline Dataset Generation and ML Model Training, (2) Real-Time Git Commit Processing and Predictive Invalidation, and (3) Production RAG Integration and Vector Store Refresh. 

This multi-phase architectural design intentionally decouples computationally expensive model training tasks from live commit processing. By reserving model training for offline repository mining, the system executes real-time commit invalidation in sub-0.15 seconds per commit pair, making it ideal for background integration into developer CI/CD workflows and production vector database serving backends.

The complete system workflow is illustrated in the System Architecture Flowchart (Fig. \ref{fig:architecture}), while the class component interaction contracts are detailed in Fig. \ref{fig:class_diagram}.

\begin{figure}[htbp]
\centering
\begin{tikzpicture}[node distance=0.9cm and 1.2cm, auto, >=stealth',
    phase/.style={rectangle, draw, fill=blue!15, text width=3.3cm, text centered, rounded corners, minimum height=0.7cm, font=\bfseries\tiny},
    block/.style={rectangle, draw, fill=gray!10, text width=3.1cm, text centered, rounded corners, minimum height=0.6cm, font=\tiny},
    cloud/.style={ellipse, draw, fill=green!15, text width=2.2cm, text centered, minimum height=0.6cm, font=\tiny}]
   
    % Nodes
    \node [phase] (p1) {PHASE 1: Offline Training\\(Commit Sampling, AST Graph, BFS Features, RF Regressor)};
    \node [phase, below=1.0cm of p1] (p2) {PHASE 2: Live Commit Invalidation\\(Cosmetic Filter, Fast BFS, ML Predict, BFS Decay)};
    \node [phase, below=1.0cm of p2] (p3) {PHASE 3: Production Refresh\\(LlamaIndex Pipeline, PredictiveCacheFilter)};
    \node [cloud, right=1.2cm of p3] (vectordb) {Vector Database\\(ChromaDB / Qdrant)};

    % Paths
    \draw [->, thick, double] (p1) -- node[left, font=\tiny] {Model Weights} (p2);
    \draw [->, thick, double] (p2) -- node[left, font=\tiny] {Stale Node List} (p3);
    \draw [->, thick, double] (p3) -- node[above, font=\tiny] {Selective Update} (vectordb);
\end{tikzpicture}
\caption{End-to-End 3-Phase Pipeline Architecture for Predictive Semantic Cache Invalidation.}
\label{fig:architecture}
\end{figure}

\begin{figure}[htbp]
\centering
\begin{tikzpicture}[node distance=0.8cm and 1.0cm, auto, >=stealth',
    comp/.style={rectangle, draw, fill=yellow!15, text width=3.4cm, text centered, rounded corners, minimum height=0.7cm, font=\tiny}]
    
    \node [comp] (parser) {\textbf{RepoParser (Pass 1 \& 2)}\\+entities: Dict [AST]\\+graph: nx.DiGraph\\+parse\_directory()};
    \node [comp, below=0.7cm of parser] (extractor) {\textbf{FeatureExtractor (BFS Engine)}\\+pagerank\_impact: Dict\\+extract\_features()\\+\_get\_code\_metrics()};
    \node [comp, below=0.7cm of extractor] (predictor) {\textbf{DriftPredictor (ML + BFS Decay)}\\+model: RandomForest\\+predict\_drift()\\+weighted\_bfs\_decay()};
    \node [comp, below=0.7cm of predictor] (plugin) {\textbf{PredictiveCacheFilter (LlamaIndex)}\\+stale\_entities: Set\\+\_\_call\_\_()};

    \draw [->, thick, dashed] (parser) -- (extractor);
    \draw [->, thick, dashed] (extractor) -- (predictor);
    \draw [->, thick, dashed] (predictor) -- (plugin);
\end{tikzpicture}
\caption{Component Interaction Class Diagram across Pipeline Phases.}
\label{fig:class_diagram}
\end{figure}

\subsection{Phase 1: Offline Dataset Generation and ML Model Training}
Phase 1 establishes the predictive foundation of our framework by mining historical repository commit pairs to train a machine learning model capable of accurately predicting vector space drift.

\subsubsection{Historical Repository Ingestion and Commit Pair Sampling}
The pipeline ingests a target software repository (e.g., \texttt{psf/black}) and samples historical Git commit pairs $(C_i \to C_{i+1})$ across fixed commit strides. By analyzing real developer commit histories over time, the system captures natural patterns of code churn, function refactorings, parameter additions, and logic modifications that occur during typical software evolution.

\subsubsection{AST Call Graph Construction and Two-Pass Symbol Resolution}
Python's native \texttt{ast} parser parses every source file in the repository to extract fine-grained code entities (individual functions and class methods). A two-pass symbol resolution procedure (\texttt{RepoParser}) converts these raw ASTs into a directed call graph \texttt{nx.DiGraph}:
\begin{itemize}
    \item **Pass 1 (Entity Discovery):** Walks the AST to extract function names, line boundaries, parameter lists, return annotations, and docstrings, compiling a repository-wide entity symbol table.
    \item **Pass 2 (Call Edge Resolution):** Inspects call expressions within each function body. It maps call targets against imported module symbol tables and matches dotted method paths and suffixes (\texttt{\_resolve\_call\_target}), constructing directed caller-to-callee edges.
\end{itemize}

\subsubsection{BFS Graph Feature Extraction Engine (Role \#1 of BFS)}
Starting from Git-modified root nodes, a **Breadth-First Search (BFS) graph traversal engine** traverses the call graph to compute structural proximity and impact metrics for every function in the repository. BFS algorithms compute shortest-path hop distances (\texttt{distance\_to\_modified\_directed}, \texttt{distance\_to\_modified\_undirected}), inspect 1-hop BFS neighbor modifications (\texttt{modified\_dependents\_count}, \texttt{modified\_dependencies\_count}), and evaluate random-walk reachability via NetworkX Personalized PageRank (\texttt{pagerank\_impact}) seeded specifically on edited nodes.

\subsubsection{Ground-Truth Drift Labeling and Random Forest Regressor Training}
Neural code embedding models (\texttt{microsoft/unixcoder-base}, \texttt{all-MiniLM-L6-v2}) generate dense vector embeddings for each entity before commit $C_i$ ($v_f^{(i)}$) and after commit $C_{i+1}$ ($v_f^{(i+1)}$). The ground-truth target label is computed as the continuous cosine distance drift:
\begin{equation}
D_{\text{cosine}}(v_f^{(i)}, v_f^{(i+1)}) = 1 - \frac{v_f^{(i)} \cdot v_f^{(i+1)}}{\|v_f^{(i)}\|_2 \|v_f^{(i+1)}\|_2}
\end{equation}
A Random Forest regressor ($N=100$ decision trees) is trained on the 25 extracted features to learn how structural code edits predict vector distance drift, and its trained model weights are serialized to disk (\texttt{predictor.py}).

\subsection{Phase 2: Real-Time Git Commit Processing and Predictive Invalidation}
When a software developer pushes a new Git commit, Phase 2 executes live, sub-second predictive cache invalidation in five sequential sub-stages:

\subsubsection{Stage 1: AST Cosmetic Normalization and Filtering}
Modified source files undergo round-trip AST compilation. Stage 1 strips comments, docstrings, and formatting whitespace. If a function's compiled AST structure is identical to its previous state, non-semantic cosmetic edits are filtered out immediately, preventing unnecessary re-embeddings.

\subsubsection{Stage 2: Sub-0.15s Fast Feature Computation}
For all affected entities, Stage 2 calculates the 25 static graph and topological features (including BFS shortest paths and PageRank impact) on the fly in **sub-0.15 seconds** per commit pair—completely eliminating the need to generate expensive neural vector embeddings during inference.

\subsubsection{Stage 3: Machine Learning Drift Inference}
The calculated feature vectors are passed through the pre-trained Random Forest regressor (\texttt{DriftPredictor}). The model infers predicted continuous cosine distance drift values $\hat{D}_{\text{cosine}}$ for every entity. Entities predicted to drift beyond threshold $\tau$ are selected for cache invalidation.

\subsubsection{Stage 4: Weighted BFS Decay Cascade Propagation (Role \#2 of BFS)}
To account for indirect semantic drift (where unedited caller functions undergo vector drift because a callee's interface changed under context-aware chunking), Stage 4 executes a Breadth-First Search decay cascade. Starting from Git-modified root nodes, invalidation signals propagate backward along reverse call-graph edges (from callees to callers). To ensure that live invalidation executes in sub-0.15 seconds without invoking neural embedding models during inference, the propagation score $S(v)$ decays at each hop according to lightweight static AST diff similarity $\text{ASTSim}(u_{\text{old}}, u_{\text{new}})$ or ML-predicted similarity $(1 - \hat{D}_{\text{cosine}}(u))$:
\begin{equation}
S(v) = S(u) \cdot \Big( \alpha \cdot \text{PPR}(v) + \beta \cdot \text{CallFreqNorm}(v, u) + \gamma \cdot \text{ASTSim}(u_{\text{old}}, u_{\text{new}}) \Big)
\end{equation}
where $\text{ASTSim}(u_{\text{old}}, u_{\text{new}})$ measures round-trip AST edit similarity without generating any neural embeddings. If decayed score $S(v) \ge \tau_{\text{decay}}$, caller $v$ is flagged for re-embedding.

\subsubsection{Stage 5: Minimal Stale Entity List Compilation}
Aggregates the directly modified entities and indirectly decayed caller entities into a single, minimal list of stale function IDs.

\subsection{Phase 3: Production RAG Integration and Vector Store Refresh}
Phase 3 operationalizes the invalidation output within enterprise RAG ingestion pipelines:

\subsubsection{LlamaIndex Pipeline Interception}
In production RAG systems, codebase updates pass through LlamaIndex using custom extensions (\texttt{llama\_index\_integration.py}):
\begin{itemize}
    \item \texttt{CodeGraphNodeParser}: Subclasses \texttt{BaseNodeParser} to parse Python codebases and attach call-graph metadata, PageRank scores, and entity IDs to \texttt{TextNode} objects.
    \item \texttt{PredictiveCacheFilter}: Subclasses \texttt{TransformComponent} to intercept code chunks inside an \texttt{IngestionPipeline}.
\end{itemize}

\subsubsection{Selective Vector Store Refresh Mechanics}
The \texttt{PredictiveCacheFilter} checks each code chunk against Phase 2's stale entity list:
\begin{itemize}
    \item **Stale Entities:** Passed to the neural embedding model and updated in ChromaDB or Qdrant.
    \item **Fresh Entities (Unchanged):** Bypasses re-embedding entirely, reusing existing cached vectors.
\end{itemize}

\subsubsection{Enterprise Operational \& Cost Advantages}
Reduces re-embedding token and compute costs by **over 90\%**, maintains sub-second decision turnaround, and guarantees 100\% vector index freshness for LLM context retrieval.

% ===============================================================
\section{Experimental Protocol and Benchmarking Architecture}
% ===============================================================

\subsection{Target Repository and Commit Sampling Protocol}
The experimental evaluation is designed around open-source Python repositories, using \texttt{psf/black} (~350 Python files) as the initial benchmark dataset. Stride sampling across Git commit pairs extracts temporal codebase snapshots to evaluate drift prediction.

\subsection{Evaluation Metrics and Comparison Baselines}
The framework will be evaluated against four comparison strategies:
\begin{enumerate}
    \item \textbf{Full Re-embed (Ground Truth):} Re-indexes 100\% of the codebase.
    \item \textbf{No Propagation (\texttt{changed\_only}):} Re-indexes directly modified files.
    \item \textbf{Hardcoded Graph Heuristic:} Re-indexes nodes based on fixed hop distance.
    \item \textbf{Learned Model (Predictive ML):} Our proposed Random Forest + BFS Decay strategy.
\end{enumerate}

\subsection{Benchmarking Pipeline Execution Architecture}
To rigorously evaluate invalidation quality and search precision, we construct a dedicated automated benchmarking pipeline (\texttt{src/benchmarking/}) comprising five modular components, as illustrated in Fig. \ref{fig:benchmarking_arch}.

\begin{figure}[htbp]
\centering
\begin{tikzpicture}[node distance=0.8cm and 1.1cm, auto, >=stealth',
    box/.style={rectangle, draw, fill=orange!15, text width=3.3cm, text centered, rounded corners, minimum height=0.65cm, font=\tiny},
    runner/.style={rectangle, draw, fill=purple!15, text width=3.3cm, text centered, rounded corners, minimum height=0.65cm, font=\tiny},
    metrics/.style={rectangle, draw, fill=green!15, text width=3.3cm, text centered, rounded corners, minimum height=0.65cm, font=\tiny}]
   
    % Nodes
    \node [box] (sampler) {\textbf{Commit Sampler Module}\\(\texttt{commit\_sampler.py})\\Stride Sampling Across Git Snapshots};
    \node [box, below=0.6cm of sampler] (query) {\textbf{Balanced Query Generator}\\(\texttt{query\_generator.py})\\20 Queries: 15 Drifted + 5 Fresh};
    \node [runner, below=0.6cm of query] (runner) {\textbf{Strategy Execution Runner}\\(\texttt{runner.py} \& \texttt{evaluator.py})\\Runs 4 Strategies vs. Mock Vector DB};
    \node [metrics, below=0.6cm of runner] (metrics) {\textbf{Metrics Calculator Module}\\(\texttt{metrics.py})\\Cost Ratio, Recall@K, Precision, Latency};

    % Paths
    \draw [->, thick] (sampler) -- (query);
    \draw [->, thick] (query) -- (runner);
    \draw [->, thick] (runner) -- (metrics);
\end{tikzpicture}
\caption{Execution Flow Architecture of Automated Benchmarking Pipeline.}
\label{fig:benchmarking_arch}
\end{figure}

\subsubsection{Commit Sampler Module (\texttt{commit\_sampler.py})}
Extracts temporal codebase snapshots ($C_i \to C_{i+1}$) across fixed commit strides (\texttt{num\_commits * commit\_stride}). It checks out each commit snapshot into isolated workspaces to evaluate feature extraction and embedding drift under realistic development histories.

\subsubsection{Balanced Query Generator (\texttt{query\_generator.py})}
Generates 20 balanced evaluation search queries for each commit pair to prevent evaluation bias:
\begin{itemize}
    \item **15 Drifted Queries:** Queries targeting functions whose ground-truth vector embeddings underwent significant drift ($D_{\text{cosine}} > \tau$).
    \item **5 Fresh Queries:** Control queries targeting unedited, non-drifted functions to verify that invalidation does not degrade search precision on fresh code.
\end{itemize}

\subsubsection{Strategy Execution Runner (\texttt{runner.py} \& \texttt{evaluator.py})}
Executes all four invalidation strategies (\texttt{full\_reembed}, \texttt{changed\_only}, \texttt{fixed\_hop}, and \texttt{predictive\_ml}) against the generated query benchmarks. It updates mock vector stores per strategy and records top-$K$ search retrieval outputs.

\subsubsection{Metrics Calculator Module (\texttt{metrics.py})}
Quantifies strategy performance across four standardized benchmark dimensions:
\begin{itemize}
    \item \textbf{Cost Ratio (\%):} $\frac{\text{Re-embedded Entities}}{\text{Total Entities in Repo}} \times 100\%$ (Measuring compute and API cost savings).
    \item \textbf{Retrieval Recall@K \& Precision@K:} Evaluating top-$K$ search retrieval accuracy against ground-truth full re-indexing.
    \item \textbf{Fidelity Error / MAE:} Mean Absolute Error between predicted drift $\hat{D}_{\text{predicted}}$ and ground-truth cosine distance $D_{\text{cosine}}$.
    \item \textbf{Decision Latency:} Per-commit execution time (in seconds) required to compute invalidation candidate lists.
\end{itemize}

% ===============================================================
\section{Risk Management, Governance, and Conclusion}
% ===============================================================

\subsection{Structured Technical Risk Management}
Potential technical risks associated with call graph construction, feature extraction, and drift prediction are identified alongside mitigation strategies in Table \ref{tab:risk_matrix}.

\begin{table}[htbp]
\caption{Structured Risk Mitigation Matrix}
\label{tab:risk_matrix}
\centering
\begin{tabular}{p{2.2cm}p{1.8cm}p{3.2cm}}
\toprule
\textbf{Identified Risk} & \textbf{Risk Impact} & \textbf{Mitigation Strategy} \\
\midrule
Incomplete Call Resolution & Missed caller dependencies & Two-pass symbol resolution matching imported modules and method suffixes. \\
False Invalidation & Extra embedding costs & Round-trip AST compilation stripping comments and docstring edits. \\
Missed Vector Drift & Retrieval quality decay & Weighted BFS Decay propagating decay scores along reverse call edges. \\
Large Repo Overhead & Processing latency & Fast static graph feature extraction in sub-0.15s without vector generation. \\
Model Prediction Error & Inaccurate invalidation & Dynamic percentile thresholding ($\tau$) and hyperparameter tuning. \\
\bottomrule
\end{tabular}
\end{table}

\subsection{Ethical Considerations and Responsible Governance}
The development and evaluation of this project adhere to ethical software engineering standards:
\begin{itemize}
    \item \textbf{Open-Source Data Usage:} All datasets and commit histories are derived strictly from public open-source software repositories.
    \item \textbf{Software License Compliance:} Repository processing respects open-source licensing terms (e.g., MIT, Apache 2.0).
    \item \textbf{Local Execution \& Privacy:} Static analysis and embedding inference execute entirely on local hardware, ensuring proprietary source code is never exposed to external cloud endpoints.
    \item \textbf{Responsible Code AI Usage:} Acknowledges inherent model and static analysis limitations, providing confidence-calibrated invalidation thresholds.
\end{itemize}

\subsection{Conclusion and Future Roadmap}
In conclusion, this project proposes a native, graph-aware predictive cache invalidation framework for vector-based code search and RAG systems. By combining AST cosmetic comment filtering, 25 topological NetworkX signals, Random Forest ML drift regression, and Weighted BFS Decay propagation, the framework addresses the critical bottleneck of index staleness. Future work will explore multi-language parsing via Tree-sitter.

\begin{thebibliography}{25}

\bibitem{Fluri2025}
B.~Fluri, M.~W{\"u}rsch, M.~Pinzger, and H.~Gall, ``A retrospective of ChangeDistiller: Tree differencing for fine-grained source code change extraction,'' \emph{IEEE Trans. Software Eng.}, vol.~51, no.~3, pp. 852--857, 2025.

\bibitem{Zhang2024}
R.~Zhang, K.~Wang, R.~Duan, and L.~Li, ``A code difference extraction method based on node feature enhancement,'' in \emph{Proc. IEEE AUTEEE}, 2024, pp. 110--115.

\bibitem{Wetzlmaier2014}
T.~Wetzlmaier, C.~Klammer, and R.~Ramler, ``Extracting dependencies from software changes: An industry experience report,'' in \emph{Proc. IEEE IWSM-Mensura}, 2014, pp. 163--168.

\bibitem{Zaixing2023}
Z.~Zhang, L.~Liu, J.~Chang, L.~Wang, and L.~Liao, ``Commit classification via diff-code GCN based on system dependency graph,'' in \emph{Proc. IEEE QRS}, 2023, pp. 476--487.

\bibitem{Jin2026}
Y.~Jin et al., ``UntCC: Untangling composite commits using structural and semantic information,'' \emph{IEEE Trans. Software Eng.}, vol.~52, no.~8, pp. 2282--2302, 2026.

\bibitem{Wu2025}
Y.~Wu, W.~Suo, S.~Feng, C.~Wu, D.~Zou, and H.~Jin, ``Fine-grained code clone detection by keywords-based connection of program dependency graph,'' \emph{IEEE Trans. Reliability}, vol.~74, no.~3, pp. 3427--3441, 2025.

\bibitem{Zhang2020}
F.~Zhang, S.-C. Khoo, and X.~Su, ``Improving maintenance-consistency prediction during code clone creation,'' \emph{IEEE Access}, vol.~8, pp. 82085--82099, 2020.

\bibitem{guo2022unixcoder}
D.~Guo et al., ``UniXcoder: Unified cross-modal pre-training for code representation,'' in \emph{Proc. ACL}, 2022, pp. 7212--7225.

\bibitem{feng2020codebert}
Z.~Feng et al., ``CodeBERT: A pre-trained model for programming and natural languages,'' in \emph{Proc. EMNLP}, 2020, pp. 1536--1547.

\bibitem{guo2020graphcodebert}
D.~Guo et al., ``GraphCodeBERT: Pre-training code representations with data flow,'' in \emph{Proc. ICLR}, 2021.

\bibitem{Ngo2023}
L.~H. Ngo, V.~Sekar, E.~Leclercq, and J.~Rivalan, ``Exploring code2vec and ASTminer for Python code embeddings,'' in \emph{Proc. IEEE SEAI}, 2023, pp. 53--57.

\bibitem{Yulong2026}
Y.~He, N.~Verbin, and S.~Kovalchuk, ``Social life of code: Modeling evolution through code embedding and opinion dynamics,'' \emph{Journal of Computational Science}, vol.~96, p. 102824, 2026.

\bibitem{Qilong2026}
Q.~Wu et al., ``Dual-stage method with memory-augmented embedding learning and attention-guided re-ranking...'' \emph{Knowledge-Based Systems}, vol.~340, p. 115677, 2026.

\bibitem{lewis2020retrieval}
P.~Lewis et al., ``Retrieval-augmented generation for knowledge-intensive NLP tasks,'' in \emph{Proc. NeurIPS}, vol.~33, 2020, pp. 9459--9474.

\bibitem{gao2023retrieval}
Y.~Gao et al., ``Retrieval-augmented generation for large language models: A survey,'' \emph{arXiv preprint arXiv:2312.10997}, 2023.

\bibitem{Ersoy2025}
P.~Ersoy and M.~Er{\c{s}}ah{\i}n, ``A comparative evaluation of RAG architectures for cross-domain LLM applications,'' \emph{IEEE Access}, vol.~13, pp. 194185--194196, 2025.

\bibitem{Mijic2026}
I.~Miji{\'c} and B.~Isakovi{\'c}, ``Design and implementation of a real-time RAG-based customer relationship management system...'' in \emph{Proc. IEEE IT}, 2026, pp. 1--4.

\bibitem{Essam2026}
M.~Essam et al., ``Trust-aware multi-agent traceability: Confidence-calibrated knowledge graphs...'' in \emph{Proc. IEEE MECO}, 2026, pp. 1--8.

\bibitem{Dang2026}
S.~Dang, C.~Chen, K.~Wu, Z.~Liu, and Y.~Yang, ``CacheSense: Freshness-aware semantic caching with selective invalidation for LLM-serving backends,'' in \emph{Proc. IEEE CAIBDA}, 2026, pp. 1522--1527.

\bibitem{Satpathy2024}
A.~Satpathy et al., ``GAMap: A genetic algorithm-based effective virtual data center re-embedding strategy,'' \emph{IEEE Trans. Green Commun. Netw.}, vol.~8, no.~2, pp. 791--801, 2024.

\bibitem{yan2024athena}
Y.~Yan et al., ``Enhancing code understanding for impact analysis by combining transformers and program dependence graphs,'' in \emph{Proc. ACM Softw. Eng. (FSE)}, 2024.

\end{thebibliography}

\end{document}
